\chapter*{CAPÍTULO 5. Evaluación de la agilidad del modelo SocialScrum mediante la técnica de juicio de expertos}
\addcontentsline{toc}{chapter}{CAPÍTULO 5. Evaluación de la agilidad del modelo SocialScrum mediante la técnica de juicio de expertos}
\setcounter{chapter}{5}
\setcounter{section}{0}
\setcounter{subsection}{0}
\setcounter{subsubsection}{0}
\renewcommand{\thesection}{\thechapter.\arabic{section}}
\renewcommand{\thesubsection}{\thesection.\arabic{subsection}}
\renewcommand{\thesubsubsection}{\thesubsection.\arabic{subsubsection}}

Este capítulo presenta el diseño de validación para el marco SocialScrum. Dado el alcance de esta investigación de pregrado, se propone un plan de validación mediante juicio de expertos que permita evaluar la coherencia, pertinencia y viabilidad del marco ágil antes de su implementación empírica.

\section{Estrategia de validación}

SocialScrum es un marco metodológico conceptual que requiere validación en múltiples niveles antes de considerarse listo para adopción generalizada:

\begin{table}[H]
    \centering
    \small
    \begin{tabular}{|l|p{5cm}|p{4cm}|c|}
        \hline
        \textbf{Nivel}          & \textbf{Objetivo}                                  & \textbf{Método}                         & \textbf{Estado} \\
        \hline
        1. Coherencia teórica   & Verificar que el marco es internamente consistente & Revisión de literatura, análisis lógico & Completado      \\
        \hline
        2. Juicio de expertos   & Evaluar claridad, pertinencia y suficiencia        & Panel de expertos (Delphi simplificado) & Diseñado        \\
        \hline
        3. Viabilidad operativa & Verificar que puede implementarse                  & Caso de estudio simulado                & Completado      \\
        \hline
        4. Efectividad empírica & Medir impacto en equipos reales                    & Estudios de caso múltiples              & Trabajo futuro  \\
        \hline
    \end{tabular}
    \caption{Niveles de validación del marco SocialScrum}
    \label{tab:niveles-validacion}
\end{table}

Los niveles 1 y 3 fueron abordados en capítulos anteriores. Este capítulo presenta el diseño del nivel 2 (juicio de expertos).

\section{Diseño de validación por juicio de expertos}

\subsection{Objetivo}

Evaluar la \textbf{claridad}, \textbf{pertinencia} y \textbf{suficiencia} de los componentes principales de SocialScrum mediante la opinión estructurada de expertos en las áreas relevantes.

\subsection{Perfil de expertos requerido}

Se requiere un panel de 6-10 expertos con al menos uno de los siguientes perfiles:

\begin{table}[H]
    \centering
    \small
    \begin{tabular}{|l|p{5cm}|c|}
        \hline
        \textbf{Perfil}           & \textbf{Requisitos mínimos}                               & \textbf{Cantidad} \\
        \hline
        Ingeniería de Software    & 1+ años en desarrollo ágil, certificación Scrum           & 2-3               \\
        \hline
        Community Smells          & Publicaciones en el área o experiencia aplicada           & 1-2               \\
        \hline
        Psicología Organizacional & Título en psicología, experiencia en equipos técnicos     & 1-2               \\
        \hline
        Coaching Ágil             & 3+ años como Agile Coach, experiencia en salud de equipos & 2-3               \\
        \hline
    \end{tabular}
    \caption{Perfiles de expertos para validación}
    \label{tab:perfiles-expertos}
\end{table}

\subsection{Material a evaluar}

Los expertos evaluarán los siguientes componentes del marco:

\begin{enumerate}
    \item \textbf{Rol del Social Guide}: Perfil, competencias, proceso de selección, gobernanza.
    \item \textbf{Health Score}: Dimensiones, cálculo, frecuencia de medición, umbrales de acción.
    \item \textbf{Artefactos sociales}: Social Product Backlog, Registro de Smells, Dashboard.
    \item \textbf{Eventos adaptados}: Extensiones a Daily, Planning, Review, Retrospective.
    \item \textbf{Sistema de estimación SSP}: Dimensiones, escala de Fibonacci, conversión SSP-SP.
    \item \textbf{Salvaguardas éticas}: Confidencialidad, auditoría, protocolo con HR.
\end{enumerate}

\subsection{Instrumento de evaluación}

Para cada componente, los expertos responderán a tres criterios usando escala Likert de 5 puntos:

\begin{table}[H]
    \centering
    \small
    \begin{tabular}{|l|p{8cm}|}
        \hline
        \textbf{Criterio} & \textbf{Definición}                                                                                                                \\
        \hline
        Claridad          & El componente está descrito de manera comprensible, sin ambigüedades. Un equipo podría implementarlo sin aclaraciones adicionales. \\
        \hline
        Pertinencia       & El componente es relevante para abordar la deuda social en equipos de desarrollo. Tiene sentido incluirlo en el marco.             \\
        \hline
        Suficiencia       & El componente cubre adecuadamente el aspecto que pretende abordar. No faltan elementos importantes.                                \\
        \hline
    \end{tabular}
    \caption{Criterios de evaluación por componente}
    \label{tab:criterios-evaluacion-validacion}
\end{table}

Adicionalmente, se incluyen preguntas abiertas:
\begin{itemize}
    \item ``¿Qué elementos añadiría a este componente?''
    \item ``¿Qué elementos eliminaría o simplificaría?''
    \item ``¿Qué riesgos o problemas anticipa en la implementación?''
\end{itemize}

\subsection{Procedimiento}

El proceso de validación sigue un protocolo Delphi simplificado:

\begin{enumerate}
    \item \textbf{Ronda 1 - Evaluación individual}: Cada experto recibe la documentación completa del marco y el instrumento de evaluación. Plazo: 2 semanas.

    \item \textbf{Análisis de resultados}: Se calculan promedios, desviaciones estándar y coeficiente de concordancia (Kappa de Fleiss) por componente y criterio.

    \item \textbf{Ronda 2 - Consenso}: Si hay componentes con baja concordancia ($\kappa < 0.60$) o puntuación baja (promedio $< 3.5$), se comparten los comentarios (anonimizados) y se pide reconsideración.

    \item \textbf{Revisión del marco}: Se incorporan las mejoras sugeridas por los expertos con alta concordancia.
\end{enumerate}

\subsection{Criterios de aceptación}

\begin{table}[H]
    \centering
    \small
    \begin{tabular}{|l|c|p{5cm}|}
        \hline
        \textbf{Resultado} & \textbf{Criterio}                                       & \textbf{Acción}                       \\
        \hline
        Aprobado           & Promedio $\geq$ 4.0 y $\kappa \geq$ 0.70                & El componente se mantiene             \\
        \hline
        Revisar            & 3.5 $\leq$ Promedio $<$ 4.0 o 0.60 $\leq \kappa <$ 0.70 & Revisar según comentarios de expertos \\
        \hline
        Rediseñar          & Promedio $<$ 3.5 o $\kappa <$ 0.60                      & Rediseño significativo requerido      \\
        \hline
    \end{tabular}
    \caption{Criterios de aceptación por componente}
    \label{tab:criterios-aceptacion}
\end{table}

\section{Análisis preliminar de coherencia}

Previo a la validación formal por expertos, se realizó un análisis de coherencia interna del marco. Los resultados se resumen a continuación:

\subsection{Consistencia con principios de Scrum}

\begin{table}[H]
    \centering
    \small
    \begin{tabular}{|l|c|p{6cm}|}
        \hline
        \textbf{Principio Scrum}                          & \textbf{Cumple} & \textbf{Justificación}                                                                                             \\
        \hline
        Empirismo (transparencia, inspección, adaptación) & Sí              & Health Score proporciona transparencia; eventos permiten inspección; estrategias de mitigación permiten adaptación \\
        \hline
        Auto-organización del equipo                      & Sí              & Equipo decide priorización social; Social Guide facilita, no impone                                                \\
        \hline
        Time-boxing                                       & Sí              & Extensiones a eventos se mantienen dentro de límites de tiempo                                                     \\
        \hline
        Roles definidos                                   & Sí              & Social Guide tiene responsabilidades claras, no solapa con SM o PO                                                 \\
        \hline
        Artefactos transparentes                          & Sí              & Social Product Backlog y Dashboard son visibles para todo el equipo                                                \\
        \hline
    \end{tabular}
    \caption{Verificación de consistencia con principios de Scrum}
    \label{tab:consistencia-scrum}
\end{table}

\subsection{Consistencia con literatura de \textit{community smells}}

Los \textit{community smells} abordados por SocialScrum corresponden a los identificados en la literatura académica \cite{caballero2022communitysmells}. La taxonomía de categorías (comunicación, conocimiento, relacional, organizacional) es compatible con clasificaciones existentes.

\subsection{Consistencia con Teoría de la Autodeterminación}

El Health Score, basado en los cinco valores de Scrum, tiene correspondencia indirecta con las necesidades básicas de la TAD:

\begin{itemize}
    \item \textbf{Autonomía}: Capturada parcialmente por Coraje (libertad para expresar opiniones) y Compromiso (voluntariedad).
    \item \textbf{Competencia}: Capturada parcialmente por Foco (capacidad de concentrarse en el trabajo).
    \item \textbf{Relacionalidad}: Capturada por Apertura, Respeto y parcialmente por Compromiso.
\end{itemize}

\section{Resultados esperados}

Si la validación por juicio de expertos arroja resultados positivos (promedio $\geq$ 4.0 en todos los componentes, $\kappa \geq$ 0.70), el marco estará listo para la siguiente fase: implementación piloto en equipos reales.

Los resultados de la validación permitirán:
\begin{enumerate}
    \item Identificar componentes que requieren clarificación o simplificación.
    \item Detectar elementos faltantes que los expertos consideren necesarios.
    \item Anticipar riesgos de implementación señalados por profesionales con experiencia práctica.
    \item Fortalecer la credibilidad del marco antes de proponer su adopción.
\end{enumerate}

\section{Limitaciones del diseño de validación}

\begin{enumerate}
    \item \textbf{No sustituye validación empírica}: El juicio de expertos evalúa la coherencia conceptual, no la efectividad real del marco.

    \item \textbf{Sesgo de selección}: Los expertos disponibles pueden no representar toda la diversidad de contextos donde se aplicaría SocialScrum.

    \item \textbf{Limitación geográfica}: El acceso a expertos puede estar limitado al contexto latinoamericano.

    \item \textbf{Conocimiento del marco}: Los expertos evalúan el marco ``en papel'', sin experiencia de implementación real.
\end{enumerate}

Estas limitaciones refuerzan la necesidad de complementar esta validación con estudios empíricos en equipos reales, como se propone en el capítulo de trabajos futuros.
